% Chapter 10: Starting & Refactoring Projects
\chapter{Starting and Refactoring Projects}
\label{ch:projects}

\begin{itshape}
Two situations: a blank directory for a new project,
and a 50,000-line legacy codebase with zero tests.
Both need Claude's help. Both need different protocols.
\end{itshape}

This chapter unifies two project lifecycle phases:
greenfield (starting from scratch) and brownfield (introducing Claude
to existing code). Both follow the same principle:
incremental investment, with each step independently valuable.
The testing patterns from \cref{ch:testing} and the extension mechanisms
from \cref{ch:extending} both play central roles.

\section{Greenfield: From Zero to Production}
\label{sec:greenfield}

\subsection{Day 1: Foundation}

\convergence{Three investments on day one compound into dramatically
better outcomes over the project's lifetime:}

\begin{enumerate}
  \item \textbf{Run \code{/init}} to generate a starter CLAUDE.md.
    Claude examines your project structure and proposes conventions.
  \item \textbf{Set up permissions} in \code{.claude/settings.json}.
    Deny access to \code{.env*} and \code{secrets/}.
  \item \textbf{Add one hook}: \code{PostFileWrite $\rightarrow$ auto-format}.
    This prevents style drift from the first file onward.
\end{enumerate}

\subsection{Week 1: Patterns Emerge}

\begin{enumerate}
  \item \textbf{Add testing infrastructure.}
    Configure test runner in CLAUDE.md.
  \item \textbf{Add your first custom command} for the most common workflow.
  \item \textbf{Configure PreCommit hook} $\rightarrow$ run tests.
    The single most valuable hook for new projects.
\end{enumerate}

\subsection{Month 1: Stabilization}

\begin{enumerate}
  \item \textbf{Review and trim CLAUDE.md.}
    Remove instructions Claude consistently follows without being told.
    Add instructions for rules Claude keeps violating.
  \item \textbf{Add skills} for repeated multi-step workflows.
  \item \textbf{Set up memory} for cross-session persistence.
  \item \textbf{Document architecture} as actionable instructions, not prose.
\end{enumerate}

\margintip{Most solo projects stabilize at Tier 2--3 within one month.
Resist the urge to jump to Tier 5.}

% -------------------------------------------------------------------
\section{Brownfield: Introducing Claude to Existing Code}
\label{sec:brownfield}

\marginnote[Official]{Use Plan Mode first---let Claude read and understand
the codebase without modifying anything.\sourceurl{https://docs.anthropic.com/en/docs/claude-code/overview}}

\subsection{Start with Understanding, Not Action}

\official{Before any changes, use Plan Mode (Shift+Tab) to let Claude explore.}
Ask:
\begin{itemize}
  \item ``Explain the architecture of this project.''
  \item ``What patterns does this codebase use?''
  \item ``Where are the tests? What is the coverage?''
\end{itemize}

\marginnote[Cross-Ref]{For code archaeology techniques in brownfield
projects, see \cref{ch:thinking-partner}.}

\official{Document existing conventions in CLAUDE.md, not aspirational ones.}
A CLAUDE.md that says ``all features have unit tests'' when most feature code
is untested notebook-extracted scripts will cause inconsistent behavior---some
features tested, some not.

\subsection{Incremental Adoption}

\practitioner{Adopt Claude configuration in four steps,
each independently valuable:}

\begin{enumerate}
  \item \textbf{CLAUDE.md} documenting current state (conventions as they are)
  \item \textbf{Permissions} restricting Claude to safe operations
  \item \textbf{Hooks} enforcing existing standards (not new ones)
  \item \textbf{Skills} emerging organically as patterns stabilize
\end{enumerate}

% -------------------------------------------------------------------
\section{The Incremental Refactoring Protocol}
\label{sec:refactoring-protocol}

\practitioner{Never attempt a big-bang rewrite
(this is anti-pattern \#7 from \cref{ch:antipatterns}).}
Use four-step increments, each independently shippable:

\begin{description}
  \item[1. Extract] Isolate the function or module to refactor.
    Create a clear boundary between ``changing'' and ``not changing.''

  \item[2. Test] Write characterization tests capturing current behavior.
    Claude excels here---prompt: ``Write tests that document exactly
    what this function does today, including edge cases.''

  \item[3. Harden] Refactor with tests as safety net.
    Each change verified against the characterization tests.

  \item[4. Promote] Move refactored code into production.
    Run full regression suite. Verify no downstream breakage.
\end{description}

\begin{tipbox}[Session Sizing]
Each step fits in a 25-minute focused session.
Commit after each step.
If interrupted, the project is always in a working state.
\end{tipbox}

\subsection{Characterization Tests}

\convergence{Before changing any code, have Claude write tests that capture
existing behavior.}

\begin{minted}{text}
Prompt: "Read src/features/legacy_features.py.
Write pytest tests that document its exact
current behavior:
- What DataFrames does it accept?
- What columns does it produce?
- What happens with NaN inputs?
Do NOT write tests for how it SHOULD work.
Write tests for how it DOES work."
\end{minted}

\margintip{Characterization tests are temporary---they may test buggy behavior.
That is intentional. They exist to detect \emph{changes}, not to assert \emph{correctness}.}

\subsection{Golden Master Testing for Pipelines}

For data pipelines, a golden master captures the exact output of the original
code so that refactored code can be verified against it.

\begin{enumerate}
  \item \textbf{Capture}: Run the original pipeline on sample data.
    Save the output DataFrame as a fixture (CSV or Parquet).
  \item \textbf{Compare}: New code must produce identical output
    within floating-point tolerance.
  \item \textbf{Verify}: Use \code{pd.testing.assert\_frame\_equal(new, golden, atol=1e-6)}.
\end{enumerate}

\begin{minted}{text}
Prompt: "Run legacy_build_features() on
tests/fixtures/sample_orders.csv.
Save the output DataFrame to
tests/fixtures/golden_features.parquet.
Then write a test that runs the new
build_features() on the same input and asserts
pd.testing.assert_frame_equal(
    new, golden, atol=1e-6)."
\end{minted}

\textbf{When to use}: Refactoring feature pipelines, migrating notebook code
to modules, replacing pandas with polars.

\begin{whybox}[Golden Masters Test Behavior---Not Correctness]
A golden master preserves the original output---including any bugs.
If the old code computed a feature incorrectly, the golden master encodes
that incorrect behavior. Fix bugs \emph{after} migration, with explicit tests
for the fix. The golden master's job is to guarantee that the migration
itself changed nothing.
\end{whybox}

\margintip{Use \code{atol}/\code{rtol} for floats, exact match for categoricals,
\code{check\_like=True} for column-order independence.}

% -------------------------------------------------------------------
\section{Visual: The Incremental Refactoring Cycle}

\begin{figure}[H]
\centering
\begin{tikzpicture}[
  phasebox/.style={draw=PrimaryBlue, fill=PrimaryBlue!8, rounded corners=6pt,
    minimum width=2.5cm, minimum height=1.2cm, font=\small\bfseries, align=center},
  note/.style={font=\scriptsize, text=DarkText, align=center},
  arr/.style={->, very thick, PrimaryBlue!60},
  >=Stealth
]
  \node[phasebox] (extract) at (0,2.5) {1. Extract\\[-2pt]\normalfont\scriptsize Isolate target};
  \node[phasebox] (test) at (3.5,2.5) {2. Test\\[-2pt]\normalfont\scriptsize Characterize};
  \node[phasebox] (harden) at (3.5,0) {3. Harden\\[-2pt]\normalfont\scriptsize Refactor};
  \node[phasebox] (promote) at (0,0) {4. Promote\\[-2pt]\normalfont\scriptsize Ship};

  \draw[arr] (extract) -- (test);
  \draw[arr] (test) -- (harden);
  \draw[arr] (harden) -- (promote);
  \draw[arr] (promote) -- (extract) node[midway, left, note] {Next\\module};

  \node[font=\footnotesize\itshape, text=TipGreen!80!black, align=center] at (1.75,1.25)
    {Commit after\\each step};

  \node[font=\scriptsize, text=PrimaryBlue!70] at (1.75,-0.8)
    {Each step $\approx$ 25 min $\cdot$ Always shippable};
\end{tikzpicture}
\caption{The incremental refactoring cycle. Each step is independently shippable.}
\label{fig:refactoring-cycle}
\end{figure}

% -------------------------------------------------------------------
\section{The Quality Flywheel}
\label{sec:flywheel}

\practitioner{After each refactoring sprint:}

\begin{enumerate}
  \item Update CLAUDE.md with lessons learned
  \item Add new hooks for discovered gotchas
  \item Refine skills based on what worked
  \item Update metrics baselines
\end{enumerate}

Each sprint improves not just the code but the tooling around the code.
After three sprints, the refactoring process itself is significantly faster.
For large codebases, consider parallel refactoring across independent
modules using git worktrees (\cref{sec:worktrees}).

% -------------------------------------------------------------------
\begin{keyinsight}
Refactoring with Claude follows the same principle as refactoring without it:
small, verifiable increments. The difference is that Claude can write
characterization tests in minutes instead of hours, making the
Extract $\rightarrow$ Test $\rightarrow$ Harden $\rightarrow$ Promote cycle
dramatically faster.
\end{keyinsight}

% -------------------------------------------------------------------
\begin{trybox}
Pick the messiest feature engineering function in your project (under 100 lines).
\begin{enumerate}
  \item Have Claude write characterization tests without telling it
    what the function \emph{should} do. Just: ``Document what DataFrames
    it accepts, what columns it produces, and what happens with NaN inputs.''
  \item Run the tests. Do they pass?
  \item If yes: save a golden master fixture and you have a safety net
    for refactoring. If no: you have discovered undocumented behavior.
    Either way, you have learned something valuable.
\end{enumerate}
\end{trybox}
